\chapter*{Список сокращений и условных обозначений}
\addcontentsline{toc}{chapter}{Список сокращений и условных обозначений}
\noindent
\addtocounter{table}{-1}

\begin{longtable}{@{}>{\bfseries}rl@{}}
	% Сокращения
	PDR         & Packet Delivery Rate (коэффициент доставки пакетов) \\
	QoS         & Quality of Service (качество обслуживания) \\
	IoT         & Internet of Things (интернет вещей) \\
	
	% Математические обозначения (раскомментируйте при необходимости)
	%$a_n$, $b_n$       & Коэффициенты разложения Ми в дальнем поле, соответствующие электрическим и магнитным мультиполям. \\
	%$\hat{\mathbf{e}}$ & Единичный вектор. \\
	%$E_0$              & Амплитуда падающего поля. \\
	%$j$                & Тип функции Бесселя. \\
	%$k$                & Волновой вектор падающей волны. \\
	%$L$                & Общее число слоёв. \\
	%$l$                & Номер слоя внутри стратифицированной сферы. \\
	%$\lambda$          & Длина волны электромагнитного излучения в вакууме. \\
	%$n$                & Порядок мультиполя. \\
	%$\mu$              & Магнитная проницаемость в вакууме. \\
	%$r,\theta,\phi$    & Полярные координаты. \\
	%$\omega$           & Частота падающей волны. \\
\end{longtable}
